\chapter{Einleitung}
\label{chap:einleitung}

Ein Betrieb verteilt eingehende Aufträge auf mehrere gleich schnelle Maschinen, ein Rechenzentrum
verteilt Rechenaufträge auf die Knoten eines Clusters.
In beiden Fällen ist dieselbe Frage zu beantworten: Welche Aufgabe gehört auf welche Maschine?
Entscheidend ist dabei nicht die Summe der Bearbeitungszeiten, sondern der Zeitpunkt, zu dem die
zuletzt beschäftigte Maschine ihre Arbeit beendet.
Dieser Zeitpunkt heißt \emph{Makespan}, und seine Minimierung ist der Gegenstand dieser Arbeit.

Formal sind $n$ Aufgaben mit bekannten Bearbeitungszeiten und $m$ identische Maschinen gegeben.
Jede Aufgabe läuft ohne Unterbrechung auf genau einer Maschine; gesucht ist eine Zuordnung mit
kleinstmöglichem Makespan.
In der Notation von \textcite{graham1979} trägt dieses Problem den Namen $P\,\|\,C_{\max}$.
Kapitel~\ref{chap:grundlagen} führt diese Notation zusammen mit den übrigen Bezeichnungen ein.
Das zugehörige Entscheidungsproblem ist NP-vollständig \cite{garey_johnson1979}; sofern
$\mathrm{P} \ne \mathrm{NP}$ gilt, berechnet folglich kein Algorithmus mit polynomieller Laufzeit
den minimalen Makespan.
Gefragt sind deshalb Verfahren, welche schnell eine Zuordnung liefern und deren
\emph{Approximationsgüte} sich beweisen lässt.
Diese Güte ist der Faktor, um den der erzeugte Makespan den optimalen im schlechtesten Fall
höchstens übersteigt.

\paragraph{Bekannte Ansätze.}
Der erste Algorithmus mit einer solchen Garantie stammt von \textcite{graham1966}; mit dem
Makespan-Problem nimmt die Analyse beweisbar guter Verfahren ihren Ausgang
\cite[S.~145]{hochbaum_shmoys_1987}.
Sein List-Scheduling durchläuft die Aufgaben in einer gegebenen Reihenfolge und legt jede auf
die gerade am wenigsten belastete Maschine; der erzeugte Makespan übersteigt das Optimum dabei um
höchstens den Faktor $2 - 1/m$.
Dieser Faktor ist \emph{scharf}, das heißt, kleiner lässt er sich nicht wählen: Es gibt
Instanzen, auf welchen ihm das Verhältnis von erzeugtem zu optimalem Makespan beliebig nahe
kommt \cite[S.~1571]{graham1966}.
Werden die Aufgaben zuvor absteigend sortiert, so sinkt der Faktor auf $4/3 - 1/(3m)$, und auch
dieser Wert ist scharf \cite[Theorem~2]{graham1969}.
Diese Regel ist unter dem Namen \emph{Longest Processing Time} bekannt.

Einen anderen Weg gehen \textcite{coffman1978} mit MULTIFIT.
Statt die Aufgaben unmittelbar zu verteilen, rät das Verfahren eine Schranke $C$ an den Makespan
und überlässt einem Bin-Packing-Algorithmus die Frage, ob die Aufgaben in $m$ Behälter der
Kapazität $C$ passen.
Beides ist dieselbe Frage, denn eine solche Packung ist genau eine Zuordnung, bei der keine
Maschine länger als $C$ arbeitet.
Dieser Algorithmus heißt im Folgenden \emph{Subroutine}; \textcite{coffman1978} setzen dafür
First Fit Decreasing (FFD) ein, welches die Aufgaben absteigend sortiert und jede in den ersten
Behälter legt, in den sie noch passt.
Die Packung erzeugt damit die Subroutine; eine Binärsuche über $C$, nach $k$
Halbierungsschritten abgebrochen, wählt unter den akzeptierten Schranken die kleinste aus und
gibt deren Packung zurück.
Die Güte von MULTIFIT hängt damit an einer rein kombinatorischen Größe, dem
\emph{Grenzverhältnis} $r_m$: Übersteigt eine Kapazität den optimalen Makespan um mehr als den
Faktor $r_m$, so kommt FFD mit $m$ Behältern aus (Definition~\ref{def:rm}).
Für diesen Wert weisen \textcite{coffman1978} zunächst $r_m \le 1{,}22$ nach.
Über \textcite{friesen1984} und \textcite{yue1990} ist die Schranke später auf $13/11$ verschärft
worden, und \textcite{cao1995} gibt dafür einen vollständigen Beweis;
Abschnitt~\ref{sec:chronologie} stellt diese Entwicklung im Einzelnen dar.

Beliebig nah an das Optimum führen Approximationsschemata: Zu jeder vorgegebenen Genauigkeit
$\varepsilon > 0$ liefern sie eine Zuordnung mit einem Makespan von höchstens $(1+\varepsilon)$
mal dem Optimum.
Für eine feste Maschinenzahl gibt \textcite{sahni1976} ein solches Schema an, und
\textcite{hochbaum_shmoys_1987} behandeln den Fall beliebiger Maschinenzahl über \emph{duale}
Approximation, bei der nicht die Zahl der Behälter, sondern deren Kapazität überschritten werden
darf.
Bezahlt wird die Genauigkeit mit Laufzeit: Das Schema von Hochbaum und Shmoys läuft in
$O\bigl((n/\varepsilon)^{1/\varepsilon^{2}}\bigr)$, was die Autoren selbst als nicht praktikabel
bezeichnen \cite[S.~146]{hochbaum_shmoys_1987}.
MULTIFIT liegt zwischen beiden Linien, denn seine Güte ist fest, jeder Suchschritt kostet aber
nur einen einzigen Aufruf der Subroutine.
Einmaliges Sortieren und $k+1$ solche Aufrufe ergeben eine Laufzeit in $O(k\,n \log n)$, bei
fester Suchtiefe also $O(n \log n)$ (Abschnitt~\ref{sec:multifit}).

Verbessert worden ist MULTIFIT auch, ohne die Subroutine anzutasten.
\textcite{friesen_langston1986} ersetzen den Zweig, welcher eine Kapazität verwirft, durch eine
Reparatur: Eine beschränkte Zahl besonders ungünstig belegter Maschinen wird zusammen mit den
dabei übrig gebliebenen Aufgaben optimal neu belegt.
Das Verhältnis von erzeugtem zu optimalem Makespan sinkt dadurch auf $72/61$ und liegt damit
unter $13/11$; asymptotisch ist dieser Wert scharf, und die Laufzeitordnung $O(n \log n)$
bleibt erhalten.
Diese Arbeit geht den anderen Weg und lässt die Suche unangetastet.

Welche Subroutine dort steht, ist ein eigener Strang der Literatur.
FFD benötigt höchstens $11/9$ mal so viele Behälter wie eine optimale Packung, zuzüglich einer
additiven Konstanten; diese Schranke ist scharf \cite{dosa2007}.
Für Modified First Fit Decreasing (MFFD) ist mit $71/60$ eine kleinere obere Schranke bewiesen
\cite{johnson_garey1985, yue_zhang1995}.
Da $11/9$ scharf ist, fällt das asymptotische Worst-Case-Verhältnis von MFFD damit echt kleiner
aus als jenes von FFD; ob MFFD den Wert $71/60$ überhaupt annimmt, ist dagegen offen.
Kommen in einer Instanz nur wenige verschiedene Aufgabengrößen vor, so erreicht der
$\mathrm{OPT}+1$-Algorithmus von \textcite{jansen_solisoba2010} sogar eine Packung mit höchstens
einem Behälter über dem Optimum; sein Name benennt genau diese Garantie.
Ob ein solcher Vorteil den Makespan von MULTIFIT verkleinert, ist damit allerdings nicht gesagt:
Dessen Gütegarantie beruht auf einer bestimmten Eigenschaft von FFD, welche eine schärfere
Bin-Packing-Schranke nicht ersetzt.

Daneben stehen exakte Verfahren, etwa die Arc-Flow-Modelle von \textcite{mrad2018} oder die
Suchverfahren von \textcite{akram2025}.
Sie liefern den minimalen Makespan, besitzen im schlechtesten Fall aber keine polynomielle
Laufzeit; in Kapitel~\ref{chap:evaluation} dienen ihre publizierten Werte als Vergleichsgrundlage.

\paragraph{Beitrag dieser Arbeit.}
Die Arbeit verfolgt zwei Ziele.
Das erste ist eine geschlossene Aufbereitung der Performanzanalyse von MULTIFIT.
\textcite{cao1995} dient dabei als primäre Quelle, setzt die Ergebnisse von
\textcite{coffman1978} jedoch über weite Strecken voraus und führt einzelne Beweisschritte nur
knapp aus.
Kapitel~\ref{chap:performanz} führt deshalb die Herleitung der Schranke $r_m \le 5/4$
vollständig aus und verbindet sie in Theorem~\ref{thm:guete} mit dem Makespan, den MULTIFIT nach
$k$ Suchschritten zurückgibt.
Die schärfere Schranke $13/11$ wird eingeordnet und ihre Beweisidee dargestellt, nicht aber
bewiesen; Caos Fallunterscheidung bleibt außerhalb des Umfangs dieser Arbeit.
Die zugehörige untere Schranke wird dagegen geführt, denn eine Instanz von \textcite{friesen1984}
belegt $r_m \ge 13/11$ für alle $m \ge 13$.

Das zweite Ziel betrifft die Subroutine, denn innerhalb von MULTIFIT ist sie austauschbar.
Kapitel~\ref{kap:alternative} untersucht drei Alternativen theoretisch.
Zwei Ergebnisse daraus sind eigenständige Beiträge.
Zum einen gibt das Kollaps-Lemma eine Klasse von Instanzen an, auf der MFFD und FFD innerhalb von
MULTIFIT stets dieselbe Zuordnung erzeugen.
Zum anderen wird der $\mathrm{OPT}+1$-Algorithmus von \textcite{jansen_solisoba2010} zu einem
$\varepsilon$-dualen Approximationsalgorithmus umgebaut.
Zusammen mit dem Ergebnis von \textcite{hochbaum_shmoys_1987} garantiert dieser Umbau nach $k$
Suchschritten einen Makespan von höchstens $(1+\varepsilon)(1+2^{-k})$ mal dem Optimum.
Als dritte Alternative dient das Konfigurations-LP von \textcite{gilmore_gomory1961}, welches
jede mögliche Füllung eines Behälters als eigene Variable führt und die Ganzzahligkeit zunächst
fallen lässt.
Für dieses Programm wird umgekehrt gezeigt, dass die beim Runden entstehende Schranke gerade
nicht ausreicht, um die Gütegarantie von MULTIFIT zu tragen.

Was der Austausch praktisch bewirkt, misst Kapitel~\ref{chap:evaluation} in zwei Versuchsreihen.
Die erste vergleicht FFD und MFFD auf \EvalInstanzen{} Instanzen aus fünf publizierten
Benchmark-Datensätzen, die zweite vier Subroutinen für Instanzen mit wenigen verschiedenen
Aufgabengrößen auf \HmInstanzen{} eigens erzeugten Instanzen.
Kein gemessenes Verhältnis überschreitet dabei $13/11$ - für MULTIFIT+FFD liegt dieser Wert
knapp unter der bewiesenen Schranke $13/11 + 2^{-k}$, für die übrigen Kombinationen ist er
lediglich eine Vergleichslinie.
Zugleich zeigen beide Reihen, dass aus einer schärferen Bin-Packing-Garantie kein kleinerer
Makespan folgt.

\paragraph{Aufbau.}
Kapitel~\ref{chap:grundlagen} legt die Grundlagen: Problemstellung und Notation, Bin Packing mit
FFD sowie der Algorithmus MULTIFIT samt seinem Startintervall.
Kapitel~\ref{chap:performanz} enthält die Performanzanalyse, Kapitel~\ref{kap:alternative} die
drei alternativen Subroutinen und Kapitel~\ref{chap:evaluation} die empirische Auswertung.
Kapitel~\ref{chap:fazit} fasst die Ergebnisse zusammen und benennt die offen gebliebenen Fragen.
Sämtliche Implementierungen und Messskripte liegen unter
\url{https://github.com/hempllll/MultiFit-BscThs}; den Versuchsaufbau beschreibt
Abschnitt~\ref{sec:versuchsaufbau}.
